\documentclass[10pt,a4paper]{article}
\usepackage[T1]{fontenc}
\usepackage[utf8]{inputenc}
\usepackage{mathptmx}
\usepackage{courier}
\usepackage[english]{babel}
\usepackage{amsmath,amssymb}
\usepackage[margin=1.9cm]{geometry}
\usepackage{microtype}
\usepackage{booktabs,tabularx,array}
\usepackage{titlesec}
\usepackage{caption}
\pagestyle{empty}
\titleformat{\section}{\normalsize\bfseries\scshape}{\thesection.}{0.5em}{}
\titlespacing*{\section}{0pt}{1.1ex plus .3ex}{0.5ex}
\setlength{\parskip}{0.3ex}
\renewcommand{\arraystretch}{1.08}
\newcolumntype{Y}{>{\raggedright\arraybackslash}X}
\captionsetup{font=small,labelfont=bf,skip=3pt}
\setlength{\textfloatsep}{8pt}
\setlength{\intextsep}{6pt}

\begin{document}
\begin{center}
{\large\bfseries Does Underlying Structure Predict Quantum Trainability and Dequantization?}\\[0.3ex]
{\small Charbel A. Dargham \quad\textperiodcentered\quad DALAS, Deliverable 1 (TME4): research question and dataset}
\end{center}

\section{Research question}

\emph{Can the underlying structure of a learning problem predict, before training, whether structure-preserving
quantum models of it will be trainable and resist dequantization?} Random, structure-free quantum circuits
are untrainable at scale because their gradients vanish exponentially with the number of qubits (McClean et
al.\ 2018). Building a problem's symmetry into the circuit keeps the gradients large (Larocca et al.\ 2022;
Ragone et al.\ 2024), but the same restriction can let a classical computer imitate the model, so that it
offers no quantum advantage (Cerezo et al.\ 2025). Whether a given problem lands on the useful side is today
known only after a model has been built and run. We ask whether it can be read off the problem itself, from
four kinds of structure: symmetry, locality, spectrum and data complexity. Answering it needs two
ingredients, a large and varied collection of real learning problems, and a way to extract their structure;
this deliverable builds the first and tests the second.

\section{Data collection}

The unit of observation is a \emph{problem}: one binary learning task on $n$ sites (4 to 12, the number of
qubits a model of it would use), with its labelled inputs, a fixed 70/30 TRAIN/TEST split and its label rule.
Each input $x$ is stored with an operator $M(x)$ that defines a differential equation, and the operator fixes
the problem's type (Table~\ref{tab:types}). The type depends on the mathematics, not the domain: a MaxCut
instance encodes like a graph, while a 3-SAT formula has three-site terms a graph cannot have.

\begin{table}[h]
\centering\small
\caption{Types, sources and counts. Groups are independent sources (one dataset, network or generator).}\label{tab:types}
\begin{tabularx}{\linewidth}{@{}l l Y r r r@{}}
\toprule
Type & Operator $M(x)$ & Sources and collection & Problems & Groups & Samples\\
\midrule
GRAPH & graph Laplacian & 20 real networks (Netzschleuder API: biology, social, infrastructure, technology), cut into connected pieces of $n$ nodes; real molecules (Alchemy, ZINC via TUDataset) & 345 & 22 & 100\,782\\
COVARIANCE & precision matrix & EEG motor imagery (MOABB API: BNCI2014\_001, Cho2017), one problem per subject and session; 9 sensor datasets (UEA archive) & 404 & 11 & 87\,992\\
DIAGONAL & clause-penalty Hamiltonian & SAT formulas from 11 standard generators (random $k$-SAT, XOR, NAE, planted, community, SR) & 358 & 11 & 107\,400\\
QUANTUM & Hamiltonian, ground state & 8 spin-chain models (exact ground states on both sides of a phase transition); 6 molecules (Hartree--Fock, Jordan--Wigner) & 261 & 14 & 41\,756\\
\midrule
Total & & & 1\,368 & 58 & 337\,930\\
\bottomrule
\end{tabularx}
\end{table}

Each source was listed, screened, downloaded (with a file hash) and converted by a logged pipeline. A
candidate problem was kept only if it had at least 100 samples, at least 50 distinct inputs and at most 20\%
of equal inputs carrying different labels; every rejection is recorded with its reason (for the real
networks, 892 candidates were checked and 307 kept). Labels are either given by the source (which hand was
imagined, which class a recording belongs to) or computed exactly from the input (a graph property, the side
of a phase transition); when a source gives a score, the label is $+1$ above the TRAIN median, so the TEST data
never influence the cut. Real inputs make up 749 problems (55\%), collected through the Netzschleuder,
TUDataset and MOABB APIs and the UEA archive; the SAT and spin-chain problems are generated from standard
models, because in these two domains generated instances are how problems are produced, and every problem
carries its origin.

\section{Quantitative description and variables}

The final dataset holds 1\,368 problems, 337\,930 samples and 58 groups, with $n$ from 4 to 12 and 1\,107
distinct input sets (problems sharing their inputs and differing only in their labels). Class balance lies
in $[0.2, 0.8]$ in every problem. Problems from one group are not independent, so the group, not the problem,
is the unit of evidence in every later analysis. The variables are listed in Table~\ref{tab:vars}: the
first are recorded per problem; the second are computed from the \emph{factor graph} of each operator, a
bipartite graph with one node per site and one per term, joined when the term acts on the site, and
summarised over a problem's inputs by the median and interquartile range. All are unchanged when the
operator is multiplied by a constant.

\begin{table}[h]
\centering\small
\caption{Variables.}\label{tab:vars}
\begin{tabularx}{\linewidth}{@{}l Y@{}}
\toprule
Variable & Meaning\\
\midrule
type, $n$, samples, split & operator type, number of sites, number of labelled inputs, 70/30 TRAIN/TEST\\
class balance, label rule & share of $+1$ labels; how the labels were made\\
distinct inputs, conflict rate & number of different inputs; share of equal inputs with different labels\\
source, group, origin & where the problem comes from, its independent source, real or generated\\
equation check & check method, worst residual, number of failures (0 in all 1\,368 problems)\\
\midrule
terms per site; term sizes & interactions per site; whether terms are pairwise or many-body\\
site degree spread; support density & whether a few sites dominate; how far interactions reach\\
symmetry classes and fraction & how many sites look alike under colour refinement (1: all alike, 0: none)\\
spectral entropy; participation ratio & how many Laplacian modes of the factor graph carry its weight\\
algebraic connectivity & how hard the factor graph is to cut in two\\
\bottomrule
\end{tabularx}
\end{table}

\section{Preprocessing and choices}

The raw harvest (1\,379 problems) is never modified; a script with six stated rules builds the final dataset
from it, and the raw files' checksum is identical before and after. An audit found 0 structural failures and
five problems in the content, each handled by one choice. Identical inputs on both sides of the split, in 22
problems, would leak TEST into TRAIN: these problems were re-split with identical inputs kept on one side,
rather than deduplicated, which would have changed the population, and repeated inputs are weighted so that
each counts once in TEST. One problem below the class-balance floor was excluded. In 24\% of spin-chain
samples the ground state was degenerate, so the stored state was arbitrary and the input not reproducible:
the spin chains were re-harvested with ground states fixed inside a symmetry sector, which cut such samples
from 642 to 10 on a test run. Two molecule geometries whose Hartree--Fock calculation did not converge had
untrustworthy labels and were removed. Finally, 475 sensor samples whose operator collapses to the identity
carry no information and act as a label shortcut; they are flagged rather than removed, and every result is
reported with and without them. The final dataset passes the audit with 0 failures and 0 leaks. One known
gap remains: 10 periodic XY spin-chain problems were lost to a solver failure, since fixed, and are to be
re-harvested.

\end{document}
